\documentclass[../../main.tex]{subfiles}
\begin{document}

{\bf [解答]}

\begin{figure}[htb]
  \centering
  \begin{tikzpicture}[>=stealth, every node/.style={font=\small}]
    % ----------------------------------------------------
    % Panel 1 (左): (f(\alpha), f(\beta)) の概要図
    % ----------------------------------------------------
    \begin{scope}[shift={(0, 0)}]
      % 領域ハッチング（除外される直線以外の全平面を象徴的に示す領域）
      \fill[pattern=north east lines, pattern color=blue!25] (-2.2, -1.9) rectangle (2.2, 1.9);

      % 座標軸
      \draw[->, thick, gray!80!black] (-2.4, 0) -- (2.6, 0) node[right] {$p$};
      \draw[->, thick, gray!80!black] (0, -2.1) -- (0, 2.3) node[above] {$q$};

      % 除外される直線 q = \frac{\beta-1}{\alpha-1} p （a=0 に対応）
      \def\mSlope{0.65}
      \draw[line width=2.8pt, white] (-2.2, {-2.2*\mSlope}) -- (2.2, {2.2*\mSlope});
      \draw[line width=1.1pt, dashed, red!85!black] (-2.2, {-2.2*\mSlope}) -- (2.2, {2.2*\mSlope})
      node[pos=0.75, below right=1pt, font=\scriptsize, fill=white, inner sep=1pt] {$q = \frac{\beta-1}{\alpha-1}p$};

      % 原点記号（第2象限に配置して直線との重なりを回避）
      \node[above left=2pt] at (0, 0) {$O$};

      % 代表点 A(p, q)
      \coordinate (A) at (1.1, 1.25);
      \draw[densely dotted, thick, gray!70!black] (A) -- (1.1, 0) node[below=1pt, font=\footnotesize] {$p$};
      \draw[densely dotted, thick, gray!70!black] (A) -- (0, 1.25) node[left=1pt, font=\footnotesize] {$q$};
      \node[circle, fill=blue!85!black, inner sep=1.4pt] at (A) {};
      \node[above=2pt, blue!85!black, font=\footnotesize] at (A) {$A(f(\alpha), f(\beta))$};

      \node[below, font=\footnotesize] at (0, -2.1) {$(f(\alpha), f(\beta))$ の動く範囲（$\alpha\neq\beta$）};
    \end{scope}

    % ----------------------------------------------------
    % Panel 2 (右): y = f(x) のグラフ
    % ----------------------------------------------------
    \begin{scope}[shift={(6.0, 0)}]
      % 座標軸
      \draw[->, thick, gray!80!black] (-1.8, 0) -- (2.3, 0) node[right] {$x$};
      \draw[->, thick, gray!80!black] (0, -1.9) -- (0, 2.2) node[above] {$y$};
      \node[above right=2pt] at (0, 0) {$O$};

      % f(x) = 0.75*(x-1)*(x+1.2)
      \draw[line width=1.2pt, blue!85!black, domain=-1.6:1.95, samples=70, smooth]
      plot (\x, {0.75*(\x - 1.0)*(\x + 1.2)})
      node[right, font=\footnotesize] {$y=f(x)$};

      % 定点 (1, 0)
      \coordinate (Fix) at (1.0, 0);
      \node[circle, fill=gray!80!black, inner sep=1.3pt] at (Fix) {};
      \node[below=2pt, font=\footnotesize] at (Fix) {$1$};

      % 代表値 \alpha = -0.9, \beta = 1.6
      \def\alphaVal{-0.9}
      \pgfmathsetmacro{\fAlpha}{0.75*(\alphaVal - 1.0)*(\alphaVal + 1.2)}
      \coordinate (Palpha) at (\alphaVal, \fAlpha);

      \def\betaVal{1.6}
      \pgfmathsetmacro{\fBeta}{0.75*(\betaVal - 1.0)*(\betaVal + 1.2)}
      \coordinate (Pbeta) at (\betaVal, \fBeta);

      % 投影点線とラベル
      \draw[densely dotted, thick, gray!70!black] (\alphaVal, 0) node[above=1pt, font=\footnotesize] {$\alpha$}
      -- (Palpha) -- (0, \fAlpha);
      \node[above left=1pt, font=\footnotesize] at (0, \fAlpha) {$f(\alpha)$};
      \draw[densely dotted, thick, gray!70!black] (\betaVal, 0) node[below=1pt, font=\footnotesize] {$\beta$}
      -- (Pbeta) -- (0, \fBeta) node[left=2pt, font=\footnotesize] {$f(\beta)$};

      \node[circle, fill=blue!85!black, inner sep=1.3pt] at (Palpha) {};
      \node[circle, fill=blue!85!black, inner sep=1.3pt] at (Pbeta) {};

      \node[below, font=\footnotesize] at (0.2, -2.1) {$f(1)=0$ を満たす $y=f(x)$};
    \end{scope}
  \end{tikzpicture}
  \caption{$(f(\alpha), f(\beta))$ の動く範囲（左：破線を除く全平面）と定点 $(1,0)$ を通る $y=f(x)$ の概形（右）}
  \label{fig:locus_and_parabola}
\end{figure}

$p=f(\alpha), q=f(\beta)$とし，$A(p,q)$とおけば，$A$の動く範囲を求めればよい．
$c=-a-b$だから与えられた関数$f(x)$は
\begin{align}
  f(x)=a(x^2-1)+b(x-1)=(x-1)\{a(x+1)+b\} \quad (a\neq 0) \label{eq:1}
\end{align}
である．ただし題意より$f(x)$が二次式だから，$a \neq 0$である．
従って
\begin{align}
  \begin{cases}
    p = f(\alpha) = (\alpha-1)((\alpha+1)a + b) \\
    q = f(\beta) = (\beta-1)((\beta+1)a + b)
  \end{cases}
\end{align}
である．このとりうる値は$\alpha,\beta=1$かどうかで変化する．
$\alpha$について考えると，
\begin{align}
  p &=
  \begin{cases}
    (\alpha-1)((\alpha+1)a + b) & (\alpha \neq 1)\\
    0 & (\alpha=1)
  \end{cases}
\end{align}
である．$\alpha \neq 0$の時は$p$は$\alpha$の一次以下の式であり，
$a,b$が任意の値を動く時任意の実数をとる．
対称性から$\beta$についても同様だから，$\alpha, \beta$が$1$か否かで場合分けする．

\subparagraph{(i)$\alpha=\beta=1$の時}

$p=q=0$ゆえ$A(0,0)$である．

\subparagraph{(ii)$\alpha=1,\beta\not=1$の時}

$p=0$，$q$は任意の実数をとるので，$A$は直線$x=0$上を動く．

\subparagraph{(iii)$\alpha\not=1,\beta=1$の時}

(ii)と同様に$q=0$，$p$は任意の実数をとるので，$A$は直線$y=0$上を動く．

\subparagraph{(iv)$\alpha\not=1,\beta\not=1$の時}

$p=(\alpha-1)\{a(\alpha+1)+b\}$を$b$についてといて，
\begin{align}
  b=\dfrac{p}{\alpha-1}-a(\alpha+1) \tag{$\because \alpha\not=1$}
\end{align}
である．これを$q=(\beta-1)\{a(\beta+1)+b\}$に代入して
\begin{align}
  q&=(\beta-1)\left\{a(\beta+1)+\dfrac{p}{\alpha-1}-a(\alpha+1)\right\} \\
  &=\dfrac{\beta-1}{\alpha-1}p+(\beta-1)(\beta-\alpha)a
\end{align}
である．$a$が$a\neq 0$で任意の値を動く時の$(p,q)$の範囲を考えれば良い．
これは$\alpha=\beta$なら$q=p$となり，$\alpha\not=\beta$なら$p\not=\dfrac{\beta-1}{\alpha-1}q$であるから，
\begin{align}
  \begin{cases}
    \alpha=\beta\text{の時}&y=x \\
    \alpha\not=\beta\text{の時}&y\neq \dfrac{\beta-1}{\alpha-1}q
  \end{cases}
\end{align}
となる．
\casesend

以上をまとめて，
\begin{align}
  \begin{cases}
    \alpha=\beta=1\text{の時}&(0,0) \\
    \alpha=1\land\beta\not=1\text{の時}&x=0 \\
    \alpha\not=1\land\beta=1\text{の時}&y=0 \\
    \alpha\not=1\land\beta\not=1\land\alpha=\beta\text{の時}&y=x \\
    \alpha\not=1\land\beta\not=1\land\alpha\not=\beta\text{の時}&y\neq \dfrac{\beta-1}{\alpha-1}q
  \end{cases}
\end{align}
である．$\cdots$(答)
\newpage
\end{document}
